\section{Mathematical and Computational Methodology}
\label{sec:methodology}

This section states the actual model and computational protocol used in the
report. The preceding review explains why a reduced model is useful and where
its assumptions enter; the present section fixes the selected variables,
closure choices, solver-coordinate convention, boundary contract, numerical
realization, comparison operator, metrics, and provenance record before any
numerical outcomes are interpreted.

\subsection{Selected 1D Forward Model}
\label{subsec:selected-1d-forward-model}
\label{subsubsec:cross-sectional-reduction-selected-1d-model}

The model hierarchy identifies the 1D tier as the selected reduced forward map
for the numerical study. This subsection now records that selected model after
the hierarchy has established the 3D, multiscale, 2D, 1D, and 0D roles. The 1D
model reduces the spatial dimension by assuming that the vessel has a
distinguished axial direction and that pressure and velocity can be averaged
over cross-sections. This keeps pulse-wave and pressure-drop structure while
discarding secondary flow, separation, recirculation, and resolved wall traction.
The reduction follows standard compliant-artery and dimension-reduced modeling
arguments
\cite[Sec.~2, pp.~253--257]{FormaggiaEtAl2003OneDimensionalBloodFlow}
\cite[Ch.~3, Secs.~3.2 and~3.6]{KopplHelmig2023DimensionReducedModeling}
\cite[Sec.~3, pp.~24--26]{ShiEtAl2011BloodFlowModelReview}.

\subsubsection{Cross-Sectional Variables}

The reduced variables $A$, $Q$, $\bar u$, and $p_{\mathrm g}$ should be read as
cross-sectional outputs of the continuum fields under additional averaging,
geometry, wall-law, and profile assumptions, not as independent replacements for
the blood-dynamics setting.

\begin{figure}[!htb]
    \centering
    \captionsetup{aboveskip=2pt,belowskip=2pt}
    \figtikz{compliant-vessel}
    \caption{Compliant-vessel notation for the 1D reduction. The current
    configuration $\OmgBt$ is represented by a straightened averaging domain:
    section averages over $S(z,t)$ define $A(z,t)$, $Q(z,t)$, $\bar u(z,t)$,
    and $p_{\mathrm g}(z,t)$, while $R(z,t)$ records the radius entering the
    wall and geometry closures. The figure is a bookkeeping reduction from
    continuum fields to retained 1D variables, not a resolved wall-traction or
    secondary-flow model.}
    \label{fig:compliant-vessel-reduction}
\end{figure}

\begin{assumption}[Straight single-vessel geometry]
    \label{ass:1d-straight-vessel-geometry}
    The selected 1D model assumes a vessel of length $L$ aligned with the
    $z$-axis and a circular cross-section with radius $R(z,t)$. For
    $0<z<L$, $S(z,t)$ is the interior cross-section used for averaging, not
    a boundary component of the open lumen.
\end{assumption}

\begin{definition}[Section-averaged 1D fields]
    \label{def:section-averaged-1d-fields}
    The retained reduced variables are
    \begin{flalign*}
        \quad A(z,t) &\defined \int_{S(z,t)} 1 \, dS, &&
            \text{cross-sectional area}, \\
        \quad Q(z,t) &\defined \int_{S(z,t)} u_z \, dS, &&
            \text{volumetric flow rate}, \\
        \quad \bar u(z,t) &\defined \frac{Q(z,t)}{A(z,t)}, &&
            \text{mean axial velocity}, \\
        \quad p_{\mathrm g}(z,t) &\defined \frac{1}{A(z,t)}
            \int_{S(z,t)} p_{3D,\mathrm g} \, dS, &&
            \text{mean gauge pressure}.
    \end{flalign*}
    Here $p_{3D,\mathrm g}$ is continuum pressure in the wall-law gauge
    frame. The pressure-reference convention for reduced ratios is fixed
    separately in Convention~\ref{conv:appendix-pressure-reference-ratio}.
\end{definition}

\subsubsection{Closed Area--Flow Model and Wall Law}

\begin{definition}[Classical no-slip 1D baseline]
    \label{def:classical-no-slip-1d-baseline}
    The classical no-slip 1D baseline used as a model-family member starts from
    the fixed-wall continuum trace $\vu=0$ on the wall boundary
    $\Gamma_w\subset\pOmgB$, together with separate inlet and outlet data on
    $\Gamma_{\mathrm{in}}$ and $\Gamma_{\mathrm{out}}$. After cross-sectional
    averaging, the no-slip antecedent is represented by the parabolic
    profile/friction convention $\alpha_{\mathcal P}=4/3$ and
    $g_{\mathcal P}=4$. In the implementation manifests this member is recorded
    as \texttt{classical-1d-no-slip}; it keeps the same reduced state variables
    and selected Canic--Koiter pressure--area wall law as the extended 1D
    framework, but omits the Canic effective-alpha variable-radius correction.
    The descriptor \texttt{canic-extended-1d} records the current
    extended-model member.
\end{definition}

\begin{definition}[Closed 1D area-flow model]
    \label{def:closed-1d-area-flow-model}
    For $\mathcal W=\mathcal W_{\mathrm{elastic1D}}$ and an incompressible
    fluid with density $\rho$, the working compliant single-vessel model is
    written with the composite axial pressure derivative
    \begin{flalign*}
        \quad
        \frac{d}{dz}p_{\mathrm g}(A(z,t),z,t)
        =
        \pd_A p_{\mathrm g}(A,z,t) \pd_z A
        +
        \left.\pd_z p_{\mathrm g}(A,z,t)\right|_A. &&
    \end{flalign*}
    The reduced balance law is
    \begin{flalign*}
        \quad \pd_t A
        + \pd_z Q
        &= 0, \numberthis\label{eq:1d-mass-aq} && \\
        \quad \pd_t Q
        + \pd_z
            \left(\alpha \frac{Q^2}{A}\right)
        + \frac{A}{\rho}\frac{d}{dz}p_{\mathrm g}(A(z,t),z,t)
        &= - C_f(z,A,Q) \frac{Q}{A}.
        \numberthis\label{eq:1d-momentum-aq} &&
    \end{flalign*}
    The parameter $\alpha$ is the momentum-flux correction factor induced by
    the selected cross-sectional velocity profile, $C_f$ is the wall-friction
    closure, and $p_{\mathrm g}(A,z,t)$ is the wall-law gauge pressure
    closing the reduced system. This is the reduced
    $\mathcal W_{\mathrm{elastic1D}}$ member of the wall-model axis: wall
    response is represented through the pressure--area law rather than through a
    separate structural PDE
    \cite[Sec.~2, pp.~253--255]{FormaggiaEtAl2003OneDimensionalBloodFlow}
    \cite[Ch.~3, Sec.~3.2, pp.~40--47]{KopplHelmig2023DimensionReducedModeling}.
\end{definition}

\begin{assumption}[Velocity-profile and rheology closure]
    \label{ass:1d-first-run-alpha-friction}
    A closed 1D area-flow model must declare the momentum-flux correction
    factor $\alpha$, the wall-friction closure $C_f$, any rheology descriptor
    $\mathcal R$, and the pressure--area wall closure $\mathcal W$. These
    choices determine how cross-sectional profile effects, wall friction,
    viscosity or effective viscosity, and wall compliance enter the reduced
    equations. Numerical-record details such as solver coordinates,
    shear-rate proxies, and discrete source terms are deferred to
    Appendix~\ref{app:numerical-methods-details}.
\end{assumption}

\begin{definition}[Selected Canic--Koiter pressure--area wall law]
    \label{def:1d-elastic-wall-law}
    The selected elastic wall closure is the Canic--Koiter thin-wall
    pressure--area law. In conventional area notation it may be written
    \begin{flalign*}
        \quad p_{\mathrm g}(A,z,t)
        =
        p_{\mathrm{ext}}(z,t)
        + \frac{\beta(z)}{A_0(z)}
            \left(\sqrt{A(z,t)}-\sqrt{A_0(z)}\right).
        \numberthis\label{eq:1d-wall-law} &&
    \end{flalign*}
    Here $A_0(z)$ is the reference area, $p_{\mathrm{ext}}$ is the external
    pressure, and $\beta(z)$ is a wall-stiffness parameter. Some
    implementations absorb $A_0^{-1}$ into $\beta$; the essential closure
    is that $p_{\mathrm g}$ is known once $A_0$, $\beta$, and
    $p_{\mathrm{ext}}$ are supplied. In the radius-squared coordinate
    $a=R^2$ used by the implemented solver, the selected member is
    \begin{flalign*}
        \quad
        p_{1,\mathrm g}(a,z,t)
        =
        p_{\mathrm{ext}}(z,t)
        +
        \frac{K}{R_0(z)^2}
        \left(\sqrt a-R_0(z)\right),
        \qquad
        K=\frac{Eh}{1-\sigma^2}. &&
    \end{flalign*}
    The diagnostic extended pressure then adds the Canic variable-radius
    correction
    \begin{flalign*}
        \quad
        p_{2,\mathrm g}(a,Q,z,t)
        =
        g_{\mathcal P}\rho\nu_{\mathrm{eff}}^{\mathcal R}
        \frac{Q}{a}\frac{R_0'(z)}{R_0(z)}. &&
    \end{flalign*}
    This $p_2$ term is a variable-radius viscous correction, not a replacement
    wall law. Resolved FSI would replace the algebraic pressure--area closure
    with structural unknowns and fluid--structure interface conditions.
\end{definition}

\subsubsection{Implemented Coordinates and Stenosis Geometry}

\begin{definition}[Physical-to-solver map for the implemented 1D state]
    \label{def:physical-to-solver-map}
    The continuum 1D model uses physical area $A_{\mathrm{phys}}$ and
    physical flow $Q_{\mathrm{phys}}$. The implemented radius-squared solver
    instead stores
    \begin{flalign*}
        \quad
        a=R^2,\qquad q=\frac{Q_{\mathrm{phys}}}{\pi}.
        &&
    \end{flalign*}
    Thus
    \begin{flalign*}
        \quad
        A_{\mathrm{phys}}=\pi a,\qquad
        Q_{\mathrm{phys}}=\pi q,\qquad
        \bar u=\frac{Q_{\mathrm{phys}}}{A_{\mathrm{phys}}}=\frac{q}{a}.
        &&
    \end{flalign*}
    With $K=Eh/(1-\sigma^2)$ and the Canic conservative-form convention
    $R_0^*=R_{\max}$, the selected wall law gives the solver elastic flux
    contribution recorded in Appendix~\ref{app:num-current-solver-operator}:
    \begin{flalign*}
        \quad
        \Psi_h(a)=\frac{K}{3\rho R_{\max}^2}a^{3/2},
        \qquad
        \pd_a\Psi_h(a)=\frac{K}{2\rho R_{\max}^2}\sqrt a.
        &&
    \end{flalign*}
    This map is the convention used by the Section~\ref{sec:resolved-3d-comparison}
    velocity and flow comparisons.
\end{definition}

\begin{definition}[Stenosis reference geometry]
    \label{def:1d-stenosis-reference-geometry}
    A stenosis is encoded through the reference radius or reference area. Let
    $R_{\mathrm{base}}(z)$ be the nominal healthy radius and $s(z)\in[0,1)$
    a fractional radius-reduction profile. Then
    \begin{flalign*}
        \quad
        R_0(z) &= R_{\mathrm{base}}(z)\big(1-s(z)\big),
        &
        A_0(z) &= \pi R_0(z)^2.
        \numberthis\label{eq:1d-stenosis-area} &&
    \end{flalign*}
    The baseline idealized stenotic-vessel profile is the smooth asymmetric
    kernel
    \begin{flalign*}
        \quad
        g(z) &= z-z_s+\kappa
        \exp\left(-\frac{(z-z_a)^2}{2\sigma_a^2}\right),
        &
        \psi(z) &= \exp\left(-\lambda g(z)^4\right),
        \numberthis\label{eq:1d-stenosis-kernel} && \\
        \quad
        s(z) &= s_{\max}\psi(z), \qquad 0\leq s_{\max}<1.
        \numberthis\label{eq:1d-stenosis-profile} &&
    \end{flalign*}
    In the default centimeter-scale simulation record, $z_a=2.5\,\mathrm{cm}$,
    $z_s=3.4\,\mathrm{cm}$, $\sigma_a=1\,\mathrm{cm}$,
    $\kappa=0.95\,\mathrm{cm}$, and
    $\lambda=50\,\mathrm{cm}^{-4}$.
    This geometry is the baseline idealized simulation geometry in this report,
    not a patient-specific plaque morphology or clinical calibration.
\end{definition}

\begin{remark}[Stenosis-profile and variable-radius boundary]
    \label{rmk:1d-stenosis-profile-scope}
    The profile in Eq.~\ref{eq:1d-stenosis-profile} is $C^\infty$ and
    rapidly localized rather than compactly supported. This smooth idealized
    geometry is the baseline for the reduced simulations in this report and is
    also useful for fixed-wall 3D studies: the same analytic radius profile
    supplies controlled derivatives, wall normals, severity variation, and
    repeatable mesh-refinement families without introducing segmentation noise
    or patient-specific morphology. Canic et al.\ show that standard 1D
    variable-geometry models may miss stenotic pressure and velocity trends
    unless additional variable-radius terms are included
    \cite[Abstract; Secs.~1--2]{CanicEtAl2024Extended1DStenosis}. Those
    corrections are included only by the \texttt{canic-extended-1d}
    forward-model descriptor; \texttt{classical-1d-no-slip} retains the same
    selected wall law but omits the Canic variable-radius corrections.
\end{remark}

\subsubsection{Conservative Form and Boundary Contract}

\begin{definition}[Conservative flux/source notation for the reduced wall law]
    \label{def:1d-conservative-flux-source-notation}
    For a reduced wall closure $\mathcal W$ admitting an elastic potential,
    define
    \begin{flalign*}
        \quad
        \mathbf U
        &=
        (A,Q)^\top, &&
        \mathbf F_{\mathcal W}(\mathbf U,z,t)
        =
        \begin{pmatrix}
            Q\\[2pt]
            \alpha Q^2/A+\Psi_{\mathcal W}(A,z,t)
        \end{pmatrix}, &&\\
        \quad
        \mathbf S_{\mathcal W,\mathcal R}(\mathbf U,z,t)
        &=
        \begin{pmatrix}
            0\\[2pt]
            S_{\mathrm{geom}}^{\mathcal W}(A,z,t)
            +S_f^{\mathcal R}(A,Q,z)
        \end{pmatrix}. &&
    \end{flalign*}
    The potential and sources are tied to the selected closure descriptors by
    \begin{flalign*}
        \quad
        \pd_A \Psi_{\mathcal W}(A,z,t)
        &=
        \frac{A}{\rho} \pd_A p_{\mathrm g}^{\mathcal W}(A,z,t), &&\\
        \quad
        S_{\mathrm{geom}}^{\mathcal W}(A,z,t)
        &=
        \left.\pd_z \Psi_{\mathcal W}(A,z,t)\right|_A
        -
        \frac{A}{\rho}
        \left.\pd_z p_{\mathrm g}^{\mathcal W}(A,z,t)\right|_A,
        &&\\
        \quad
        S_f^{\mathcal R}(A,Q,z)
        &=
        -C_f^{\mathcal R}(z,A,Q)\frac{Q}{A}. &&
    \end{flalign*}
    For $\mathcal W_{\mathrm{elastic1D}}$,
    $p_{\mathrm g}^{\mathcal W}$ is the wall law in
    Definition~\ref{def:1d-elastic-wall-law}. In component form,
    \begin{flalign*}
        \quad \pd_t A + \pd_z Q &= 0,
        \numberthis\label{eq:1d-conservative-mass} && \\
        \quad \pd_t Q
        + \pd_z \left(
            \alpha\frac{Q^2}{A}+\Psi_{\mathcal W}(A,z,t)
        \right)
        &=
        S_{\mathrm{geom}}^{\mathcal W}(A,z,t)
        +S_f^{\mathcal R}(A,Q,z),
        \numberthis\label{eq:1d-conservative-momentum} &&
    \end{flalign*}
    The wall descriptor $\mathcal W$ determines the pressure--area law,
    elastic potential, wave speed, and geometry source; the rheology descriptor
    $\mathcal R$ determines the viscosity or friction closure. A resolved FSI
    model is therefore not obtained by merely relabeling
    $S_{\mathrm{geom}}$: it introduces structural state and interface laws.
    The displayed split records the balance-law objects needed by finite-volume
    schemes, while a computed realization may still use a particular
    well-balanced or source-split discretization
    \cite{LeVeque2002FiniteVolumeMethodsForHyperbolicProblems}
    \cite[Summary; Secs.~2.1--2.3; Secs.~5--6]{BoileauEtAl2015Benchmark}
    \cite{FormaggiaEtAl2007FSICoupling}.
\end{definition}

\begin{definition}[Single-vessel boundary data]
    \label{def:1d-single-vessel-boundary-data}
    For a single-vessel solve on $z\in[0,L]$, prescribe one inlet condition
    and one outlet condition:
    \begin{flalign*}
        \quad Q(0,t) &= Q_{\mathrm{in}}(t)
        \quad \text{or} \quad
        \bar u(0,t) = \bar u_{\mathrm{in}}(t), && \\
        \quad p_{\mathrm g}(L,t) &= p_{\mathrm{out},\mathrm g}(t)
        \quad \text{or} \quad
        p_{\mathrm g}(L,t) = p_{\mathrm{WK,g}}(Q(L,\cdot),t). &&
    \end{flalign*}
    The second outlet option denotes a terminal Windkessel-style relation. The
    numerical study uses prescribed $Q_{\mathrm{in}}(t)$ and
    $p_{\mathrm{out},\mathrm g}(t)$ because they give an unambiguous reduced
    single-vessel contract.
\end{definition}

\begin{definition}[Characteristic speeds for the reduced elastic subsystem]
    \label{def:1d-characteristic-speeds}
    Let $\lambda_\pm$ denote the two characteristic speeds of the
    positive-area elastic 1D subsystem obtained from the homogeneous
    $(A,Q)$ balance law. Their formula and assumptions are given in
    Appendix~\ref{app:num-elastic-potential-form}.
\end{definition}

\begin{assumption}[Subcritical boundary regime]
    \label{ass:1d-subcritical-boundary-regime}
    The baseline boundary treatment is used only in the subcritical regime
    $\lambda_-<0<\lambda_+$, where one characteristic enters through the inlet
    and one characteristic enters through the outlet.
\end{assumption}

Under this regime, a finite-volume scheme constrains the incoming information at
each boundary and leaves the outgoing state determined by the interior solve.
This is the reduced-model analogue of declaring enough boundary data for the
selected equation class; it is not, by itself, a validation of a particular
boundary discretization.


\subsection{Source-to-Implementation Model Variant}
\label{subsec:source-to-implementation-variant}

The selected implementation is a documented Canic-derived extended model, not
a literal reproduction of every source-paper parameter and normalization in
\cite{CanicEtAl2024Extended1DStenosis}. The report therefore uses the phrase
\emph{$R_{\max}$-normalized Canic-derived extended model} when source fidelity
matters. Table~\ref{tab:source-to-implementation-model} records the relevant
source-to-implementation choices.

\begin{table}[!htb]
    \centering
    \scriptsize
    \caption{Source-to-implementation map for the selected 1D model.}
    \label{tab:source-to-implementation-model}
    \begin{tabularx}{\textwidth}{@{}p{0.24\textwidth}p{0.34\textwidth}X@{}}
        \toprule
        \textbf{Feature} & \textbf{Canic et al. source form} & \textbf{Implemented report realization} \\
        \midrule
        Pressure denominator
            & Elastic pressure written with a local $R_0(z)^{-2}$ denominator.
            & Evolution wall potential and source use the constant
            $R_{\max}^{-2}$ normalization specified in
            Appendix~\ref{app:num-current-solver-operator}. \\
        Profile parameters
            & Source simulations state $\gamma=9$ and $\alpha=1.1$.
            & Principal runs use the parabolic profile
            $\alpha_{\mathcal P}=4/3$, $g_{\mathcal P}=4$; power-profile runs
            are sensitivity/diagnostic records. \\
        Friction coefficient
            & Profile-dependent wall/friction closure in the reduced model.
            & Implemented through the declared velocity-profile shear factor
            and rheology descriptor. \\
        $p_2$ differentiation
            & Variable-radius pressure correction differentiated in the model.
            & Locally frozen-viscosity realization for the $p_2$ derivative. \\
        External pressure
            & External pressure term may be included in the formulation.
            & $p_{\mathrm{ext}}\equiv0$ in the reported runs. \\
        Coordinates
            & Physical area and flow notation in the source paper.
            & Solver state is $(a,q)$ with
            $A_{\mathrm{phys}}=\pi a$, $Q_{\mathrm{phys}}=\pi q$. \\
        Boundary rule
            & Boundary conditions are model/problem dependent.
            & Fixed-area characteristic approximation applied with persisted
            subcritical-regime diagnostics. \\
        Discrete well balancing
            & Continuous rest equilibrium is a model-level property.
            & Rusanov/MUSCL is not exactly well balanced for spatially varying
            $R_0$; Chapter~\ref{sec:numerical-verification} reports
            zero-forcing rest-state drift. \\
        \bottomrule
    \end{tabularx}
\end{table}

\subsection{Selected and Supported Discretizations}
\label{subsec:muscl-ssprk3-discretization}

The reported C23/C40 comparison uses the finite-volume realization specified in
Appendix~\ref{app:num-current-solver-operator}. In summary, the radius-squared
state $(a,q)$ is advanced on a one-dimensional grid with a MUSCL reconstruction,
minmod limiting, source splitting for the wall/geometry and variable-radius
terms, and native SSPRK3 time integration. The numerical comparison runs use
$N=400$ cells on $0\le z\le6$ cm, a timestep cap of $10^{-5}$ s, and a CFL cap
of 0.45. The boundary-state approximation is the fixed-area characteristic
contract recorded in Appendix~\ref{app:num-boundary-state-realization}.

The local Julia project also exposes the report-now spatial and time-integration
surface in Table~\ref{tab:implemented-discretization-surface}. These options are
implementation-check and sensitivity context unless a result table explicitly
identifies them as the active method. The project dependencies include
\texttt{SciMLBase} and \texttt{OrdinaryDiffEq}; they do not include an external
WENO or ENO finite-volume package, so WENO/ENO reconstructions are future
method-development options rather than report-now methods.

\begin{table}[!htb]
    \centering
    \scriptsize
    \caption{Implemented report-now discretization surface. The principal
    comparison in this report uses MUSCL finite volume with minmod limiting and
    native SSPRK3; the other entries are benchmark, sensitivity, or backend
    parity context supported by the local Julia project.}
    \label{tab:implemented-discretization-surface}
    \begin{tabularx}{\textwidth}{@{}p{0.20\textwidth}p{0.30\textwidth}X@{}}
        \toprule
        \textbf{Surface} & \textbf{Implemented option} & \textbf{Report role} \\
        \midrule
        Spatial
            & First-order finite-volume Rusanov
            & Low-order cell-average baseline and DG $p=0$ bridge. \\
        Spatial
            & MUSCL finite volume with minmod-limited interface states and
            Rusanov flux
            & Principal C23/C40 comparison method and main convergence target. \\
        Spatial
            & Richtmyer/Lax--Wendroff finite volume with minmod-limited
            interface states
            & Native fixed-step sensitivity method; not exposed through the
            SciML \texttt{ODEProblem} backend because its flux predictor depends
            on $\Delta t/\Delta z$. \\
        Spatial
            & Modal DG with Legendre degree $p=0,1,2$
            & Native DG benchmark rows; $p=0$ reduces to the cell-average FV
            path, while $p=1,2$ are high-order spatial context rather than the
            principal comparison method. \\
        Time
            & Forward Euler, SSPRK2, and SSPRK3
            & Native fixed-step options; SSPRK3 is the selected reported
            integrator, while Euler and SSPRK2 are diagnostic sensitivity rows. \\
        Time
            & SciML/OrdinaryDiffEq \texttt{auto}, \texttt{tsit5}, and
            \texttt{rodas5p}
            & Adaptive backend-parity checks for semi-discrete operators that
            the current \texttt{ODEProblem} path supports. \\
        \bottomrule
    \end{tabularx}
\end{table}

\subsection{Cross-Sectional Velocity Reconstruction}
\label{subsec:cross-sectional-velocity-reconstruction}

The 1D state directly supplies a section mean,
$\bar u_{1D}(z,t)=q(z,t)/a(z,t)$, and physical flow
$Q_{1D}(z,t)=\pi q(z,t)$. Principal radial profile plots use the
current-radius coordinate $r/\sqrt{a(z,t)}$ and the parabolic-profile
reconstruction
\[
    u_{1D}(r;z,t)=2\,\frac{q(z,t)}{a(z,t)}
    \left(1-\left(\frac{r}{\sqrt{a(z,t)}}\right)^2\right)
\]
with the supplemental reference-radius coordinate $r/R_0(z)$ retained as an
operator-control output. That reconstruction is a profile diagnostic, not a
resolved secondary-flow, separation, skewness, or wall-traction model.

\subsection{Plane--Tetrahedron Comparison Operator}
\label{subsec:plane-tetrahedron-comparison-operator}

For a target plane $S_j=\Omega_{\mathrm{3D}}\cap\{z=z_j\}$, the default
\texttt{CrossSectionQuadratureOperator} intersects every tetrahedron with
$S_j$, linearly interpolates axial velocity on the cut edges, sorts the cut
polygon vertices in the plane, and triangulates the polygon by a centroid fan.
Let $\mathcal Q_j$ denote the resulting cut triangles. The reported 3D area,
flow, and mean velocity are
\begin{flalign*}
    \quad
    A_{3D,j}
    &=
    \sum_{\Delta\in\mathcal Q_j}|\Delta|,
    &
    Q_{3D,j}
    &=
    \sum_{\Delta\in\mathcal Q_j}|\Delta|\,\bar u_{z,\Delta},
    &
    \bar u_{3D,j}
    &=
    Q_{3D,j}/A_{3D,j}. &&
\end{flalign*}
For a cut triangle $\Delta$ whose vertices carry reconstructed axial
velocities $u_{z,1}$, $u_{z,2}$, and $u_{z,3}$,
\begin{flalign*}
    \quad
    \bar u_{z,\Delta}
    =
    \frac{u_{z,1}+u_{z,2}+u_{z,3}}{3}. &&
\end{flalign*}
The corresponding 1D quantities use physical flow, not the raw solver flow
coordinate:
\begin{flalign*}
    \quad
    Q_{1D,j}
    =
    \pi q(z_j,t),
    \qquad
    \bar u_{1D,j}
    =
    q(z_j,t)/a(z_j,t). &&
\end{flalign*}
Radial profile rows use the same cut triangles, binned by area-conservative
three-point triangle quadrature in either $r/\sqrt{a(z,t)}$ or $r/R_0(z)$.
The emitted controls include 10/20/40-bin sensitivity, area-weighted bin
means, within-bin azimuthal variance, and the current-to-reference area
mismatch. Empty bins are reported as invalid rows rather than zero-velocity
observations.

\subsection{Numerical Experiments, Discrepancy Metrics, and Provenance}
\label{subsec:numerical-experiments-metrics-provenance}

The resolved-velocity comparison evaluates the selected 1D solution against two
available resolved 3D velocity comparison datasets at nominally aligned
final-time samples: the 1D output is sampled at $1.0$ s and the XDMF velocity
field records $0.9995$ s. The cases are labelled C23 and C40 by stenosis
severity. They are computational cross-model diagnostics, not pressure-flow
results and not clinical validation.

\begin{center}
\fbox{\begin{minipage}{0.90\linewidth}
\small
\textbf{Solver-coordinate convention.} In the comparison record, the stored
solver state \texttt{A} is $a=R^2$, not physical circular area. The physical
map is Definition~\ref{def:physical-to-solver-map}:
$A_{\mathrm{phys}}=\pi a$, $Q_{\mathrm{phys}}=\pi q$, and
$\bar u=q/a=Q_{\mathrm{phys}}/A_{\mathrm{phys}}$.
\end{minipage}}
\end{center}

All comparison numbers come from
\path{simulations/output/3d_comparison/full_t1_native_nx400/}, interpreted
under Convention~\ref{conv:computed-realization-record}. The regeneration
command, source-path convention, and file hashes are listed in
Appendix~\ref{app:code-and-ai-use}. The raw XDMF/HDF5 inputs remain local
ignored files under \path{simulations/data/3d/canic_case3/}; the report
archives the generated CSV and PGFPlots-ready static data.

\begin{table}[!htb]
    \centering
    \scriptsize
    \caption{Case matching for the resolved-velocity comparison. The case
    labels C23 and C40 are used in the main text and figures; the source ID
    records the upstream local data directory.}
    \label{tab:t1-case-matching}
    \begin{tabularx}{\textwidth}{@{}llXlrr@{}}
        \toprule
        \textbf{Label}
            & \textbf{Source ID}
            & \textbf{Velocity metadata}
            & \textbf{Operator}
            & \textbf{3D time (s)}
            & \textbf{Time offset (s)} \\
        \midrule
        C23 & 77 & \path{77/velocity.xdmf}
            & \texttt{CrossSectionQuadratureOperator}
            & 0.9994999999999453 & $5.0\times 10^{-4}$ \\
        C40 & 60 & \path{60/velocity.xdmf}
            & \texttt{CrossSectionQuadratureOperator}
            & 0.9994999999999453 & $5.0\times 10^{-4}$ \\
        \bottomrule
    \end{tabularx}
\end{table}

\begin{landscape}
\begin{table}[!htb]
    \centering
    \small
    \caption{Model-matching matrix for the 1D--3D velocity diagnostic. Status
    labels record what the archived report artifact establishes, not what a
    future matched study might supply.}
    \label{tab:t1-model-matching-matrix}
    \begin{tabularx}{\linewidth}{@{}
        >{\raggedright\arraybackslash}p{0.13\linewidth}
        >{\raggedright\arraybackslash}p{0.23\linewidth}
        >{\raggedright\arraybackslash}p{0.23\linewidth}
        >{\raggedright\arraybackslash}p{0.11\linewidth}
        >{\raggedright\arraybackslash}X@{}}
        \toprule
        \textbf{Feature}
            & \textbf{1D model}
            & \textbf{3D comparison dataset}
            & \textbf{Status}
            & \textbf{Implication} \\
        \midrule
        Reference-profile geometry
            & Same analytic stenosis profile and declared severity.
            & Source mesh metadata is interpreted against the same nominal
            profile.
            & approximately matched
            & The intended stenosis shape is shared. \\
        Cut-operator area
            & Static reference area $A_{\mathrm{ref}}=\pi R_0^2$ is available.
            & Tetrahedral cuts audit against the same static reference area.
            & approximately matched
            & Cut areas show typical sub-percent agreement with localized
            multi-percent outliers. \\
        Current lumen area
            & Dynamic 1D area is $A_{1D}=\pi a(z,t)$.
            & Fixed 3D cut area is $A_{3D}(z)$.
            & unknown
            & The area denominator mismatch relevant to $\bar u=Q/A$ is not
            decomposed in the archived artifact. \\
        Wall model
            & Compliant reduced Canic--Koiter wall law.
            & Velocity field on a fixed geometry; matching wall-law metadata
            is not persisted here.
            & unmatched
            & Wall effects cannot be isolated from velocity mismatch. \\
        Density
            & $\rho$ recorded by the 1D parameters.
            & Not persisted in the tracked comparison tables.
            & unknown
            & Density-sensitive pressure or wave-speed claims are out of scope. \\
        Viscosity/rheology
            & Newtonian $\nu=0.04$ cm$^2$/s.
            & Rheology and viscosity metadata are not persisted in the
            report artifact.
            & unknown
            & The velocity comparison is not a material-parameter validation. \\
        Initial condition
            & Geometry-rest state, then forced by the inlet.
            & Prior transient history is not persisted.
            & unmatched
            & Final-time differences may include start-up history mismatch. \\
        Inlet condition
            & Reference-area flow
            $q_{\mathrm{in}}=R_0(0)^2\bar u_{\mathrm{in}}$ from nominal
            $\bar u_{\mathrm{in}}=22.5$ cm/s.
            & Boundary condition used for the 3D field is not persisted here.
            & unknown
            & Flow mismatch cannot be attributed uniquely to closure discrepancy. \\
        Outlet condition
            & Fixed-area characteristic approximation.
            & Outlet treatment is not persisted here.
            & unknown
            & Reflections and downstream constraints remain unquantified. \\
        Time / transient history
            & Evaluated at $T=1.0$ s from a start-up transient.
            & XDMF time is $0.9995$ s; earlier history is not recorded.
            & unknown
            & The final timestamps are close, but phase/history matching is not
            proved. \\
        Numerical resolution
            & $N=400$ finite-volume cells.
            & Node-centered velocity fields with loaded node counts reported.
            & unknown
            & 1D cell count and 3D node count are not comparable resolution
            measures. \\
        Velocity observable
            & Section mean $\bar u=q/a$ and physical flow $\pi q$.
            & Plane--tetrahedron cross-section quadrature of axial velocity.
            & matched
            & The compared observable is well defined. \\
        Pressure availability
            & Pressure-output vocabulary exists but is not used here.
            & No pressure comparison is reported.
            & unmatched
            & This is not a pressure-drop, FFR, or CT-FFR study. \\
        \bottomrule
    \end{tabularx}
\end{table}
\end{landscape}

\begin{table}[!htb]
    \centering
    \scriptsize
    \caption{Shared 1D realization and output-operator parameters for the
    $t=1.0$ comparison. Realized CFL, characteristic-speed, mass, and
    positivity diagnostics are persisted in the generated comparison and
    verification records.}
    \label{tab:t1-parameter-matching}
    \begin{tabularx}{\textwidth}{@{}p{0.22\textwidth}X@{}}
        \toprule
        \textbf{Quantity} & \textbf{Recorded value} \\
        \midrule
        Forward model & \path{canic-extended-1d}. \\
        Wall model & $\mathcal W_{\mathrm{elastic1D}}$ with the selected
        Canic--Koiter pressure--area law
        \path{canic-koiter-thin-membrane}. \\
        Rheology & Newtonian rheology with $\nu=0.04$ cm$^2$/s. \\
        Velocity profile & Parabolic profile with
        $\alpha_{\mathcal P}=4/3$. \\
        Variable-radius terms & Canic effective-alpha and geometry-source
        corrections, including the locally frozen $p_2$ derivative in the
        extended member. \\
        Physical parameters & $\rho=1.055$ g/cm$^3$, $E=5.02\times10^6$
        dyn/cm$^2$, $h=0.06$ cm, $\sigma=0.5$,
        $K=4.016\times10^5$ dyn/cm, $R_{\max}=0.18$ cm, and
        $p_{\mathrm{ext}}\equiv0$. \\
        Discretization & $N=400$, $L=6$ cm, MUSCL finite volume with minmod
        limiter, native SSPRK3 time integration. \\
        Time controls & $T=1.0$ s from a start-up transient, timestep cap
        $10^{-5}$ s, CFL cap 0.45. \\
        Initial and boundary state & Geometry-rest initial state
        $a=R_0^2$, $q=0$; inlet $\bar u_{\mathrm{in}}=22.5$ cm/s; outlet
        $a_{\mathrm{out}}=R_0(L)^2$ with the implemented characteristic
        approximation in Appendix~\ref{app:num-boundary-state-realization}. \\
        Section targets & 200 equally spaced axial planes on $0\le z\le 6$
        cm. \\
        Radial targets & Slices $z=1.951$, 2.451, and 2.951 cm with principal
        20-bin plots in $r/\sqrt{a(z,t)}$ plus supplemental $r/R_0(z)$ and
        10/20/40-bin sensitivity rows. \\
        Supplemental sensitivity & Node-slab arithmetic means are emitted only
        to \path{node-slab-sensitivity.csv} for half-widths 0.0075,
        0.0150753769, and 0.0301507538 cm. \\
        \bottomrule
    \end{tabularx}
\end{table}

Velocity and flow differences are reported as empirical sample metrics over
valid section targets: mean absolute discrepancy (MAE), root-mean-square
discrepancy (RMSE), and maximum absolute discrepancy. With
$d_{u,j}=\bar u_{1D,j}-\bar u_{3D,j}$, the velocity sample metrics are
\begin{flalign*}
    \quad
    \mathrm{MAE}_u
    &=
    \frac{1}{n}\sum_j |d_{u,j}|,
    &
    \mathrm{RMSE}_u
    &=
    \left(\frac{1}{n}\sum_j |d_{u,j}|^2\right)^{1/2},
    &
    \max |d_u|
    &=
    \max_j |d_{u,j}|, &&\\
    \quad
    \mathrm{rel}_{\mathrm{MAE}}
    &=
    \frac{\sum_j|d_{u,j}|}{\sum_j|\bar u_{3D,j}|},
    &
    \mathrm{rel}_{\mathrm{RMSE}}
    &=
    \frac{\left(\sum_j|d_{u,j}|^2\right)^{1/2}}
         {\left(\sum_j|\bar u_{3D,j}|^2\right)^{1/2}}. &&
\end{flalign*}
